\documentclass[11pt,a4paper]{article}
\usepackage{brainteaser}

\begin{document}
\brainteasertitle{4}{Les places dans l'avion}

\begin{statementbox}
Un avion possède $N$ sièges numérotés de $1$ à $N$, attribués aux passagers
$1,\ldots,N$. Les passagers montent dans cet ordre.

Le passager 1 a perdu son billet et choisit uniformément un siège. Chaque
passager suivant s'assoit à sa place si elle est libre; sinon, il choisit
uniformément parmi les sièges encore libres.

Quelle est la probabilité que le dernier passager trouve son propre siège
libre?
\end{statementbox}

\begin{hintsbox}
\begin{enumerate}
  \item Étudier les cas $N=2$ et $N=3$.
  \item Que se passe-t-il si le premier passager choisit le siège $1$, le
  siège $N$, ou un siège intermédiaire $k$?
  \item Après le choix de $k$, quel passager devient le nouveau passager
  désorienté?
  \item Peut-on réduire tout le processus à une course entre deux sièges
  décisifs?
\end{enumerate}
\end{hintsbox}

\correction

Le processus reste déterministe jusqu'à ce qu'un passager trouve sa place
occupée. Celui-ci choisit alors au hasard parmi les sièges libres et reprend le
rôle du premier passager.

Les seuls sièges décisifs sont:
\begin{itemize}
  \item le siège $1$: s'il est choisi en premier, tous les passagers suivants
  retrouvent leur place et le passager $N$ réussit;
  \item le siège $N$: s'il est choisi en premier, le passager $N$ échoue.
\end{itemize}

À chaque choix aléatoire pertinent, ces deux sièges encore libres jouent des
rôles symétriques. Chacun a donc la même probabilité d'être choisi avant
l'autre.

On peut aussi formaliser le raisonnement par récurrence. Si $P_N$ désigne la
probabilité cherchée,
\[
P_N=\frac1N\left(1+0+\sum_{k=2}^{N-1}P_{N-k+1}\right),
\]
et l'hypothèse $P_m=1/2$ pour les tailles inférieures redonne $P_N=1/2$.

\begin{resultbox}
\[
\boxed{\Pp(\text{le passager }N\text{ trouve sa place})=\frac12.}
\]
\end{resultbox}

\end{document}
